% Generated by `taskgen course examples` from arithmetic.montgomery_mul_even (seed 1).
% Do not edit: change examples.yaml and regenerate.
\begin{example}[метод Монтгомери для чётного модуля]\label{ex:redc-mul-even}
Вычислить произведение целых чисел по модулю: \[853194 \cdot 990419\ (\textrm{mod}\ 7600418)\] с помощью метода Монтгомери. Обосновать выбор параметра $R$.

\par\noindent\textit{Решение.} Модуль чётный: $m = 7600418 = 2^{1} \cdot 3800209$, и $R = 2^k$ не взаимно прост с $m$. Решаем сравнение по модулям $2^{1}$ и $s = 3800209$ и объединяем по КТО.\par По модулю $2^{1} = 2$ (младшие $1$ бит): $853194 \cdot 990419 \equiv 0 \cdot 1 \equiv 0\ (\textrm{mod}\ 2)$.\par По модулю $s$ методом Монтгомери (по модулю $s$: $a \equiv 853194$, $b \equiv 990419$). $R = 2^{22} = 4194304$ "--- наименьшая степень двойки, превосходящая модуль $3800209$; $R$ взаимно прост с нечётным модулем, а деление на $R$ "--- сдвиг. $m' = -m^{-1} \bmod R = 4111247$, $R^2 \bmod m = 127404$ (здесь $m = 3800209$).
\par $\mathrm{REDC}(T)$: $u = (T \bmod R)\, m' \bmod R$, $t = (T + u m)/R$, при $t \ge m$ $t \gets t - m$.
\[\begin{array}{l|r|r|r|r}
 & T & u & t & \mathrm{REDC}(T) \\ \hline
a \cdot b & 845019548286 & 845410 & 967444 & 967444 \\
x \cdot (R^2 \bmod m) & 123256235376 & 1374608 & 1274837 & 1274837 \\
\end{array}\]
Итак, $ab \equiv 1274837\ (\textrm{mod}\ 3800209)$. КТО:\[\left\{\begin{array}{l}
x \equiv 0\ (\textrm{mod}\ 2) \\
x \equiv 1274837\ (\textrm{mod}\ 3800209) \\
\end{array}\right.\qquad M = 7600418.\]
$M_{1} = M/2 = 3800209 \equiv 1$, $y_{1} = M_{1}^{-1} \equiv 1\ (\textrm{mod}\ 2)$; $M_{2} = M/3800209 = 2 \equiv 2$, $y_{2} = M_{2}^{-1} \equiv 1900105\ (\textrm{mod}\ 3800209)$.\[\begin{aligned}x &\equiv \textstyle\sum_i r_i M_i y_i \\ &= 0\cdot3800209\cdot1 + 1274837\cdot2\cdot1900105 \\ &\equiv 5075046\ (\textrm{mod}\ 7600418).\end{aligned}\]
\[853194 \cdot 990419 \equiv 5075046\ (\textrm{mod}\ 7600418).\]
\end{example}
